\chapter{Experimentation}
\label{ch:experimentation}

\section{Parameter choices and dynamics regimes}
\label{sec:params_fixed}

The experiment catalogue separates four sources of modelling difficulty.
Table~\ref{tab:main_regimes} lists the nine \emph{primary benchmark regimes}
that carry the main-text analysis in Chapters~\ref{ch:rom_results}
and~\ref{ch:wsindy_results}: five constant-speed (CS) regimes spanning
gas, blackhole, supernova, crystal, and pure-Vicsek behaviours, plus four
variable-speed (VS) diagnostic variants.
For the EF-ROM branch, these regimes test the limits of shift-aligned POD
and latent autoregression under increasing structural complexity.
For the WSINDy branch, they test whether weak-form identification can
recover the correct dominant mechanisms (diffusion, attraction, repulsion,
alignment) in each regime.
The remaining configurations in the systematic catalogue
(Appendix~\ref{app:full_regime_catalogue}) provide supplementary coverage
of force-parameter sweeps and D'Orsogna outcome-targeted families; they are
used for robustness checks but are not central to the main argument.

To keep comparisons across learning models meaningful, we fix a core set of
numerical and grid parameters across the full pipeline
(Table~\ref{tab:fixed_params}) and vary only the physics and stochastic knobs
that define each dynamics regime.  All configuration files are provided in the accompanying repository.

\begin{table}[ht]
\centering\small
\caption{Fixed parameters shared by all dynamics regimes and pipeline stages.}
\label{tab:fixed_params}
\begin{tabular}{lll}
\hline
Parameter & Value & Code key \\
\hline
Number of particles $N$                    & 300                         & \cv{N} \\
Domain size $L_x = L_y$                    & 25                & \cv{Lx}, \cv{Ly} \\
Boundary conditions                        & periodic          & \cv{bc} \\
Time step $\Delta t$                        & 0.04\,s           & \cv{dt} \\
Training horizon $T_{\mathrm{train}}$       & 20\,s (500 steps)  & \cv{T} \\
Test horizon $T_{\mathrm{test}}$            & 200\,s            & \cv{test\_sim.T} \\
Interaction radius $R$                      & 4.0               & \cv{R} \\
KDE grid resolution $n_x \times n_y$       & $64 \times 64$    & \cv{nx}, \cv{ny} \\
KDE bandwidth                               & 5.0 grid cells    & \cv{bandwidth} \\
Snapshot subsampling stride                 & 3                 & \cv{subsample} \\
Effective ROM time step $\Delta t_{\mathrm{ROM}}$ & $0.12\,$s       & $= \Delta t \times \text{subsample}$ \\
Fixed POD modes $m$                         & 19                & \cv{fixed\_modes} \\
Density transform                           & $\sqrt{\cdot}$    & \cv{density\_transform} \\
Shift alignment                             & enabled (mean ref) & \cv{shift\_align} \\
Mass post-processing                        & scale             & \cv{mass\_postprocess} \\
\hline
\end{tabular}
\end{table}

The fixed choices in Table~\ref{tab:fixed_params} are shared experimental settings
rather than per-regime optima.  They are held constant so that later
differences in ROM or WSINDy performance can be attributed primarily to
the dynamics rather than to retuning of coarse-graining, reduction, or
lifting.  In particular, the KDE bandwidth $\sigma{=}5.0$~grid cells was
selected heuristically to balance feature resolution with noise
suppression; a systematic study of the coupling between
coarse-graining scale and identified coefficients remains a subject for
future inquiry (Section~\ref{sec:future_work}).

The regime catalogue is designed to separate four distinct sources of
difficulty: coherent transport speed, stochastic disorder, strong nonlocal
clustering, and state-dependent speed modulation.  Each regime family varies
one of these axes while holding the others fixed, so that performance
differences in later chapters can be attributed to specific physical
mechanisms.

\subsection{NDYN regimes (physics-knob variations)}
\label{subsec:ndyn_regimes}

Fourteen \emph{NDYN} regimes isolate individual physics knobs---speed, noise
strength, attraction/repulsion amplitude, force length scale, and speed
mode---while holding the others at moderate baseline values.
The repository defines 14 NDYN configurations, of which 11 non-duplicate
entries are retained in the final systematic catalogue
(Appendix~\ref{app:full_regime_catalogue}).
All use Morse
forces except NDYN08, which disables forces
entirely (pure Vicsek alignment~\cite{vicsek1995vicsek}).

Pure-alignment and weak-force constant-speed regimes are expected to be the
most favourable for shift-aligned POD and compact latent autoregression,
because their density movies are dominated by a single coherent translation
that the preprocessing largely removes.  Strong-clustering and
variable-speed regimes, by contrast, are expected to stress both latent
linearity and sparse closure discovery: the former introduces near-singular
density spikes, while the latter decouples speed from heading and thereby
inflates the effective latent dimension.

\subsection{DO regimes (outcome-targeted variations)}
\label{subsec:do_regimes}

A second set of fourteen \emph{DO} (dynamics-outcome) regimes holds speed
$v_0{=}3$, noise $\eta{=}0.2$, and interaction parameters fixed except for
the repulsion coefficient~$C_r$ and repulsion length~$\ell_r$.  This
sweep targets specific collective behaviours---collective swarms, single and
double mills, double rings, and escape patterns---following the D'Orsogna
classification~\cite{dorsogna2006morse}.  The full DO catalogue is provided
in Appendix~\ref{app:full_regime_catalogue}.

\subsection{Primary benchmark regimes}
\label{subsec:main_regimes}

The main-text analysis in Chapters~\ref{ch:rom_results}
and~\ref{ch:wsindy_results} focuses on nine regimes that span the key
behavioural axes (disorder, attraction, repulsion, equilibrium, alignment only)
in both constant-speed~(CS) and variable-speed~(VS) variants.
Table~\ref{tab:main_regimes} lists these regimes with their governing
parameters; the complete systematic catalogue (11~non-duplicate NDYN
+ 14~DO, most with VS counterparts; 37 produced usable results) is provided in
Appendix~\ref{app:full_regime_catalogue}.

\begin{table}[ht]
\centering
\caption{Primary benchmark regimes used for the main-text analysis.
  Speed mode is constant-with-forces for CS variants and
  variable for VS variants.  Shared parameters:
  $\mu_t{=}0.3$, $r_{\mathrm{cut}}{=}5\max(\ell_r,\ell_a)$; see
  Table~\ref{tab:fixed_params} for the remaining fixed settings.}
\label{tab:main_regimes}
\resizebox{\textwidth}{!}{%
\begin{tabular}{llrrrrrrl}
\hline
Label & Nickname & $v_0$ & $\eta$ & $C_a$ & $C_r$ & $\ell_a$ & $\ell_r$ & Behaviour \\
\hline
NDYN04      & gas          & 3.0 & 1.5  & 0.5  & 0.2  & 1.5 & 0.5 & disordered gas \\
NDYN04\_VS  & gas (VS)     & 3.0 & 1.5  & 0.5  & 0.2  & 1.5 & 0.5 & variable-speed gas \\
NDYN05      & blackhole    & 4.0 & 0.15 & 20.0 & 0.05 & 1.5 & 0.3 & near-singular spike \\
NDYN05\_VS  & blackhole (VS)& 4.0 & 0.15 & 20.0 & 0.05 & 1.5 & 0.3 & variable-speed collapse \\
NDYN06      & supernova    & 3.0 & 0.15 & 0.05 & 15.0 & 1.5 & 0.5 & explosive repulsion \\
NDYN06\_VS  & supernova (VS)& 3.0 & 0.15 & 0.05 & 15.0 & 1.5 & 0.5 & variable-speed explosion \\
NDYN07      & crystal      & 2.0 & 0.1  & 3.0  & 3.0  & 2.0 & 0.5 & balanced equilibrium \\
NDYN07\_VS  & crystal (VS) & 2.0 & 0.1  & 3.0  & 3.0  & 2.0 & 0.5 & variable-speed equilibrium \\
NDYN08      & pure Vicsek  & 3.0 & 0.3  & ---  & ---  & --- & --- & alignment only \\
\hline
\end{tabular}}%  end resizebox
\end{table}

\noindent\textbf{Variable-speed variants.}
\label{subsec:vs_variants}
Most NDYN and DO configurations have a companion \emph{\_VS} variant
that changes exactly one parameter:
the speed mode from constant-with-forces to variable.  In
variable-speed mode, Morse forces modulate agent \emph{speed} as well as
heading, producing non-constant flux fields.  NDYN08 (pure Vicsek, no forces)
and NDYN14 (already variable-speed) have no \_VS counterpart.

%% ---------------------------------------------------------------------------
\section{Initial condition families}
\label{sec:ic_families}

To probe ROM robustness beyond a single ``typical'' microscopic state, we
generate ensembles of trajectories from four qualitatively different
\emph{initial position} distributions.  The intent is to
(i)~span a range of macroscopic organisations (homogeneous, localised,
symmetric, multi-modal) and (ii)~induce distinct transient regimes under the
\emph{same} simulator and numerical settings.  All particles evolve on the
periodic domain $\Omega=[0,L_x]\times[0,L_y]$.

\paragraph{Common velocity initialisation.}
Across all IC families we draw random headings so that differences in dynamics
are driven by the initial \emph{spatial} configuration:
\begin{equation}
  \theta_i \sim \mathcal{U}(0,2\pi),\qquad
  \mathbf{v}_i(0)=v_0(\cos\theta_i,\sin\theta_i)^\top,\qquad i=1,\dots,N,
  \label{eq:ic_vel_init}
\end{equation}
so $\|\mathbf{v}_i(0)\|=v_0$ for all particles and the initial polarisation is
close to zero in expectation.

\paragraph{Uniform IC.}
\label{subsec:ic_uniform}
The uniform family represents a maximum-entropy
baseline with no imposed spatial structure:
\begin{equation}
  \mathbf{x}_i(0)\sim \mathcal{U}(\Omega), \qquad i=1,\dots,N,
  \label{eq:ic_uniform_def}
\end{equation}
equivalently $x_i\sim\mathcal{U}(0,L_x)$ and $y_i\sim\mathcal{U}(0,L_y)$
independently.  The expected coarse density is spatially constant,
$\mathbb{E}[\rho(\mathbf{r},0)] = N/(L_xL_y)$.  Uniform ICs serve both as
training diversity and as the baseline for generalisation comparisons.

\paragraph{Single Gaussian IC.}
\label{subsec:ic_gaussian}

The single-Gaussian family produces a localised
``blob'' with controllable centre and spread:
\begin{equation}
  \mathbf{x}_i(0)\sim \mathcal{N}(\boldsymbol{\mu},\sigma^2 I_2),
  \qquad i=1,\dots,N.
  \label{eq:ic_gaussian_def}
\end{equation}
The centre $\boldsymbol{\mu}=(\mu_x,\mu_y)$ is drawn from a $4\times 4$ grid
of interior positions and
$\sigma^2$ is varied across three values, yielding
$4\times 4\times 3=48$ parameter combinations with 2 samples each,
giving 96~training runs per regime.

\paragraph{Ring IC.}
\label{subsec:ic_ring}

The ring family enforces rotational symmetry at
$t{=}0$ and typically generates vortex-like macroscopic patterns.  Let
$\mathbf{c}=(L_x/2,\,L_y/2)$ be the domain centre.  Angles and radii are
sampled in polar coordinates:
\begin{equation}
  \theta_i \sim \mathcal{U}(0,2\pi),\qquad
  r_i \sim \mathcal{N}(R_{\mathrm{ring}},\,\sigma_r^2),
  \label{eq:ic_ring_polar}
\end{equation}
and mapped to Cartesian positions
$\mathbf{x}_i(0)=\mathbf{c}+r_i(\cos\theta_i,\sin\theta_i)^\top$, then
wrapped into~$\Omega$ via modulo.  Training enumerates five radii and
four widths with 4 samples each, yielding $5\times 4\times 4=80$~training
runs per regime.

\paragraph{Two-cluster Gaussian IC.}
\label{subsec:ic_two_clusters}

The two-cluster family generates a bimodal initial
distribution as an equal-weight mixture of two isotropic Gaussians:
\begin{equation}
  p(\mathbf{x})
  =\tfrac12\,\mathcal{N}(\boldsymbol{\mu}_1,\sigma^2 I_2)
  +\tfrac12\,\mathcal{N}(\boldsymbol{\mu}_2,\sigma^2 I_2),
  \label{eq:ic_two_clusters_mix}
\end{equation}
implemented by sampling $N_1=\lfloor N/2\rfloor$ particles from the first
component and $N_2=N-N_1$ from the second.  The centres are placed
symmetrically about the domain midpoint, separated horizontally by distance~$d$:
\begin{equation}
  \boldsymbol{\mu}_1=\Big(\tfrac{L_x}{2}-\tfrac{d}{2},\,\tfrac{L_y}{2}\Big),
  \qquad
  \boldsymbol{\mu}_2=\Big(\tfrac{L_x}{2}+\tfrac{d}{2},\,\tfrac{L_y}{2}\Big).
  \label{eq:ic_two_clusters_centers}
\end{equation}
Default values are $d=L_x/3$ and $\sigma=\min(L_x,L_y)/8$.  After sampling,
positions are wrapped into~$\Omega$ via modulo.  Training enumerates
four separations and three widths with 7 samples each, yielding
$4\times 3\times 7=84$~training runs per regime.

%% ---------------------------------------------------------------------------
\noindent\textbf{Training--test split summary.}
\label{subsec:train_test_split}
Table~\ref{tab:ic_split} summarises the per-regime dataset composition.
The four held-out test runs are intended as balanced regime-level diagnostic
stress tests across IC families, not as a Monte Carlo estimate of
out-of-sample uncertainty; they probe whether each model class preserves
fidelity under qualitatively different initial spatial organisations.
Every
regime (NDYN and DO alike) uses the identical IC schedule, so the total
training set comprises $340 \times 25 = 8{,}500$ simulations across the
25~base regimes in the experiment catalogue, each with a companion variable-speed variant where applicable.

\begin{table}[htbp]
\centering
\caption{Training and test runs per dynamics regime.}
\label{tab:ic_split}
\begin{tabular}{lrr}
\hline
IC family & Training runs & Test runs \\
\hline
Gaussian     & 96 & 1 \\
Uniform      & 80 & 1 \\
Ring         & 80 & 1 \\
Two-clusters & 84 & 1 \\
\hline
\textbf{Total} & \textbf{340} & \textbf{4} \\
\hline
\end{tabular}
\end{table}

Each run is assigned a unique identifier from which a deterministic RNG seed
is derived, ensuring exact reproducibility.  The density movies produced by
these simulations feed both modeling branches.  For the EF-ROM pipeline they
are assembled into a global snapshot matrix for POD basis construction, as
described in Chapter~\ref{ch:global_pod_pipeline}.  For WSINDy, the same
regime-specific training and test simulations are retained in physical space
and converted into density and auxiliary multi-fields for weak-form regression.

These IC families expose specific modelling biases:
uniform ICs probe spurious structure formation; Gaussian ICs test localised
transport; ring ICs test rotationally organised dynamics not reducible to
translation; and two-cluster ICs test interaction between separated peaks.
Ring and two-cluster ICs are expected to be harder for the shift-aligned
POD pipeline, since translational alignment does not resolve all coherent
spatial organisation.
